\chapter{Retinal Detachment Surgery}

\section*{Introduction \& Clinical Indications}
Retinal detachment surgery is a critical ophthalmic intervention aimed at reattaching the neurosensory retina to the underlying retinal pigment epithelium (RPE) when it has become separated. This separation, termed retinal detachment, can occur due to various mechanisms, including rhegmatogenous detachment, which is caused by a break in the retina, tractional detachment resulting from abnormal fibrous membranes, or exudative detachment due to fluid accumulation without a break. The surgical procedure is vital for preserving vision, as untreated retinal detachment can lead to irreversible vision loss or blindness. The surgery's primary objective is to restore the anatomical and functional integrity of the retina, ensuring that the photoreceptors receive adequate metabolic support from the choroid. Depending on the type and severity of the detachment, surgical techniques may include scleral buckling, pneumatic retinopexy, or pars plana vitrectomy, each tailored to the specific clinical scenario.

\begin{itemize}
  \item Retinal Reattachment Surgery
  \item Scleral Buckling Procedure
  \item Pars Plana Vitrectomy for Retinal Detachment (PPV for RD)
  \item Pneumatic Retinopexy
  \item Vitrectomy with Gas or Oil Tamponade
  \item Retinal Repair Surgery
\end{itemize}

The primary surgical principle in retinal detachment surgery is to restore the apposition between the neurosensory retina and the retinal pigment epithelium, allowing the photoreceptors to regain metabolic support from the choroid. This is achieved by addressing the underlying causes of the detachment. In cases of rhegmatogenous detachment, the surgical goal is to seal all retinal breaks to prevent further ingress of vitreous fluid into the subretinal space and to relieve any vitreoretinal traction that may be pulling the retina away.

Sealing is typically accomplished through techniques such as cryotherapy or laser photocoagulation, which create chorioretinal adhesions around the breaks. Additionally, the procedure aims to reduce vitreous traction on the retina, either through external indentation with a scleral buckle or internal removal of the vitreous gel during a vitrectomy. Following the sealing of breaks and relief of traction, a tamponading agent, such as an intraocular gas bubble or silicone oil, is introduced into the vitreous cavity to hold the retina against the RPE while the chorioretinal adhesions mature. The rationale is to restore both the anatomical and physiological integrity of the retina, thereby preserving or improving visual function.

The evolution of retinal detachment surgery has been marked by significant milestones over the past century. In the 1920s, Swiss ophthalmologist Jules Gonin pioneered the treatment of retinal detachment by recognizing the importance of retinal tears and demonstrating that cauterizing these tears could lead to successful reattachment. His work laid the groundwork for modern retinal surgery.

In the 1950s, Ernst Custodis and Charles Schepens advanced the field with the development of scleral buckling techniques. Custodis introduced external scleral plombage, while Schepens refined the approach with encircling bands and sponge techniques, which provided enhanced support to the retina. The introduction of pars plana vitrectomy by Robert Machemer in the early 1970s marked another significant breakthrough, allowing direct access to the vitreous cavity for the removal of vitreous traction and internal drainage of subretinal fluid. This innovation expanded the surgical options available for complex detachments. Over the years, advancements in instrumentation, such as smaller gauge vitrectomy systems and wide-angle viewing systems, have further improved surgical precision and outcomes.

Retinal detachment surgery is indicated for several clinical scenarios, including:

\begin{itemize}
  \item \textbf{Rhegmatogenous Retinal Detachment (RRD):} The most common indication, caused by a full-thickness break in the retina allowing vitreous fluid to enter the subretinal space.
  \item \textbf{Tractional Retinal Detachment (TRD):} Often associated with proliferative diabetic retinopathy or other conditions where fibrous membranes pull the retina away from the RPE.
  \item \textbf{Combined Rhegmatogenous and Tractional Retinal Detachment:} A complex scenario requiring both traction relief and closure of retinal breaks.
  \item \textbf{Giant Retinal Tears:} Large tears involving significant portions of the retinal circumference, often requiring vitrectomy with perfluorocarbon liquids.
  \item \textbf{Retinal Dialysis:} A disinsertion of the retina from the ora serrata, often traumatic, leading to detachment.
  \item \textbf{Macula-on Retinal Detachment:} An ophthalmic emergency requiring surgery within 24-48 hours to preserve central vision.
  \item \textbf{Macula-off Retinal Detachment:} Urgent but can be performed within days to weeks; earlier repair is generally preferred.
  \item \textbf{Vitreous Hemorrhage Obscuring Retinal View:} Vitrectomy may be performed to clear hemorrhage and allow for repair.
\end{itemize}

Retinal detachment surgery is primarily a first-line treatment for most forms of retinal detachment, particularly rhegmatogenous detachments. For macula-on detachments, it is considered an emergency procedure, typically performed within 24-48 hours of diagnosis to preserve central vision. In cases of macula-off detachments, while still urgent, surgery may be performed within a week, as the urgency is slightly less immediate.

In instances of tractional retinal detachment, surgery is often the primary treatment once the traction becomes significant enough to threaten or cause detachment. However, it may also serve as a secondary or salvage procedure in cases of recurrent detachment or complex situations involving severe proliferative vitreoretinopathy (PVR). Pneumatic retinopexy, while a primary treatment, is reserved for specific uncomplicated detachments with superior retinal breaks.

\section*{Relevant Surgical Anatomy}
The retina is a thin, multi-layered neurosensory tissue that lines the posterior segment of the eye, responsible for converting light into neural signals. It is approximately 0.5 mm thick at the optic disc and thins to about 0.1 mm at the ora serrata. The retina is loosely attached to the underlying retinal pigment epithelium (RPE), which is firmly anchored to Bruch's membrane and the choroid. This anatomical relationship is crucial, as the space between the neurosensory retina and the RPE is where fluid accumulates in cases of retinal detachment.

Key anatomical structures include:

\begin{itemize}
  \item **Macula:** The central region of the retina responsible for high-acuity vision, with the fovea at its center containing the highest density of cone photoreceptors.
  \item **Optic Disc:** The point where retinal nerve fibers converge to form the optic nerve, and where the central retinal artery and vein enter and exit the eye.
  \item **Vitreous Body:** A gel-like substance filling the posterior cavity, which can exert traction on the retina, leading to detachment.
  \item **Choroid:** A vascular layer between the retina and sclera that provides metabolic support to the outer retina.
  \item **Sclera:** The tough, fibrous outer coat of the eye, providing structural integrity and serving as an anchor for extraocular muscles and scleral buckles.
  \item **Ora Serrata:** The jagged anterior edge of the retina, marking the transition to the non-photosensitive ciliary body epithelium.
\end{itemize}

Understanding these structures and their relationships is essential for the successful surgical repair of retinal detachments.

\section*{Preoperative Workup \& Preparation}
Preoperative preparation for retinal detachment surgery is essential, particularly given the urgent nature of many cases. A comprehensive ophthalmic examination, including a dilated fundus examination and B-scan ultrasonography, is performed to map the retinal detachment and identify all breaks. Optical coherence tomography (OCT) may be utilized to assess macular involvement.

Patients undergo a thorough medical evaluation, including a review of their medical history, current medications, allergies, and a physical examination. Blood tests and an electrocardiogram (ECG) may be ordered. Anesthesia clearance is obtained to assess the patient's fitness for anesthesia.

\begin{itemize}
  \item \textbf{Medication Adjustments:} Patients on anticoagulants may need adjustments or temporary discontinuation.
  \item \textbf{NPO Status:} Patients are typically instructed to be NPO for 6-8 hours before surgery.
  \item \textbf{Prophylactic Antibiotics:} Topical antibiotic drops may be prescribed before surgery.
  \item \textbf{Informed Consent:} A detailed discussion regarding the procedure, risks, benefits, and expected outcomes is conducted, and informed consent is obtained.
\end{itemize}

\textbf{Absolute Contraindications:}
\begin{itemize}
  \item Severe systemic illness or unstable medical condition that precludes safe administration of anesthesia.
  \item Patient refusal of the procedure after thorough discussion.
  \item Long-standing retinal detachment with no light perception and evidence of severe proliferative vitreoretinopathy (PVR).
  \item Active, uncontrolled ocular infection that has not been adequately treated.
\end{itemize}

\textbf{Relative Contraindications \& Precautions:}
\begin{itemize}
  \item Small, asymptomatic, very peripheral retinal detachments that may be observed.
  \item Uncontrolled glaucoma that may be exacerbated by intraocular gas or silicone oil tamponade.
  \item Significant ocular inflammation that should be controlled before surgery.
  \item Poor patient compliance, particularly for post-operative head positioning.
  \item Certain systemic conditions that may complicate positioning during surgery.
  \item Pregnancy, where the risks of anesthesia must be carefully weighed.
  \item Media opacities that severely limit surgical view.
\end{itemize}

\section*{Surgical Technique \& Procedure}
Retinal detachment surgery is performed under local anesthesia with sedation or general anesthesia, depending on the patient's health and the complexity of the case. The patient is positioned supine, and the eye is prepped and draped in a sterile fashion. The specific steps vary based on the chosen technique.

\begin{itemize}
  \item \textbf{Scleral Buckling:}
    \begin{itemize}
      \item \textbf{Anesthesia \& Exposure:} Local or general anesthesia is administered. A conjunctival peritomy is performed, and the four rectus extraocular muscles are isolated with sutures.
      \item \textbf{Break Localization:} Retinal breaks are localized using indirect ophthalmoscopy and marked on the sclera.
      \item \textbf{Chorioretinal Adhesion:} Cryotherapy is applied externally over the retinal breaks to create a chorioretinal adhesion.
      \item \textbf{Buckle Placement:} A silicone sponge or encircling band is sutured to the sclera, creating an indentation that supports the retina and closes the breaks.
      \item \textbf{Fluid Drainage (Optional):} Subretinal fluid may be drained through a small sclerotomy if necessary.
      \item \textbf{Closure:} The buckle is adjusted for appropriate tension, and the conjunctiva is closed with absorbable sutures.
    \end{itemize}
  
  \item \textbf{Pars Plana Vitrectomy (PPV):}
    \begin{itemize}
      \item \textbf{Anesthesia \& Incisions:} Local or general anesthesia is used. Three small sclerotomies are made through the pars plana for the insertion of a light pipe, infusion cannula, and vitrectomy probe.
      \item \textbf{Vitreous Removal:} The vitreous gel is removed using the vitrectomy probe.
      \item \textbf{Membrane Peeling:} Any epiretinal membranes or internal limiting membranes are carefully peeled from the retinal surface.
      \item \textbf{Fluid-Air Exchange \& Drainage:} Subretinal fluid is drained, followed by a fluid-air exchange to flatten the retina.
      \item \textbf{Laser Photocoagulation:} Laser photocoagulation is applied around all retinal breaks to create a permanent chorioretinal adhesion.
      \item \textbf{Tamponade:} A tamponading agent, such as gas or silicone oil, is injected into the vitreous cavity.
      \item \textbf{Closure:} The sclerotomies are closed, often sutureless for smaller gauges, and the conjunctiva is repositioned.
    \end{itemize}
  
  \item \textbf{Pneumatic Retinopexy:}
    \begin{itemize}
      \item \textbf{Patient Selection:} Suitable for specific uncomplicated detachments with superior retinal breaks.
      \item \textbf{Chorioretinal Adhesion:} Cryotherapy or laser photocoagulation is applied to the retinal breaks.
      \item \textbf{Gas Injection:} An intraocular gas bubble is injected into the vitreous cavity.
      \item \textbf{Head Positioning:} The patient is instructed to maintain specific head positioning for several days.
    \end{itemize}
\end{itemize}

At the conclusion of any procedure, topical antibiotics and steroids are applied, and an eye shield is placed.

\section*{Complications \& Risks}
While retinal detachment surgery is generally safe, it carries potential risks and complications, categorized as general surgical risks and procedure-specific risks.

\begin{itemize}
  \item \textbf{General Surgical Complications:}
    \begin{itemize}
      \item \textbf{Anesthesia Complications:} Adverse reactions to anesthetic agents, respiratory depression, or cardiac events.
      \item \textbf{Bleeding:} Intraocular hemorrhage, which can obscure vision or complicate surgery.
      \item \textbf{Infection:} Endophthalmitis, a rare but severe complication.
      \item \textbf{Pain:} Postoperative ocular pain, usually manageable with medication.
    \end{itemize}
  
  \item \textbf{Procedure-Specific Risks:}
    \begin{itemize}
      \item \textbf{Re-detachment:} The most common complication, often requiring further surgery.
      \item \textbf{Proliferative Vitreoretinopathy (PVR):} Formation of scar tissue that can lead to recurrent detachment.
      \item \textbf{Elevated Intraocular Pressure (IOP)/Glaucoma:} Can be transient or persistent.
      \item \textbf{Cataract Formation/Progression:} Common after vitrectomy, often requiring subsequent cataract surgery.
      \item \textbf{Diplopia (Double Vision):} More common after scleral buckling.
      \item \textbf{Macular Pucker/Epiretinal Membrane:} Formation of a membrane on the macula that can distort vision.
      \item \textbf{Hypotony:} Abnormally low intraocular pressure leading to complications.
      \item \textbf{Vision Loss:} Severe complications can lead to worse vision or complete loss.
      \item \textbf{Gas/Oil Complications:} Migration of silicone oil or emulsification leading to inflammation.
      \item \textbf{Refractive Error Changes:} Changes in prescription post-operatively.
      \item \textbf{Corneal Edema/Decompensation:} Especially in patients with pre-existing corneal disease.
    \end{itemize}
\end{itemize}

\section*{Postoperative Care \& Recovery}
Following retinal detachment surgery, patients are monitored in the post-anesthesia care unit (PACU) for recovery. Vital signs are observed, and pain management is provided as needed. The operated eye is typically covered with an eye shield.

In the first 24-48 hours, patients are usually discharged home with specific instructions, including the use of prescribed eye drops and adherence to head positioning if a gas bubble was used. Activity restrictions typically include avoiding strenuous activities and heavy lifting for several weeks. Vision may be blurry initially, especially with a gas bubble, and will gradually improve over weeks to months. Full visual recovery can take 3-6 months, and secondary procedures may be necessary.

During surgery, the surgeon documents key findings, including the type of detachment, location, size, and morphology of retinal breaks, and the presence of proliferative vitreoretinopathy (PVR). The extent of the detachment, whether it involves the macula, is also noted. While the resected retina is rarely sent for pathological examination, any peeled membranes may be analyzed to confirm the presence of glial cells and other characteristics of PVR.

A successful postoperative course is characterized by gradual resolution of symptoms and visual improvement. Patients may experience discomfort, redness, and blurry vision initially, which should be manageable with prescribed analgesics. The primary goal is anatomical reattachment of the retina, confirmed during follow-up examinations. For macula-on detachments, a return to near pre-detachment vision is often achievable, while macula-off detachments may result in some degree of permanent vision loss or distortion.

Postoperative care involves a series of follow-up visits to monitor healing and optimize outcomes.

\begin{itemize}
  \item \textbf{Frequent Postoperative Visits:} Initial follow-ups within the first few days, then weekly for the first month, and subsequently monthly.
  \item \textbf{Eye Drop Regimen:} Continued use of prescribed eye drops for several weeks, with tapering based on response.
  \item \textbf{Head Positioning and Activity Restrictions:} Adherence to head positioning and activity restrictions is crucial.
  \item \textbf{Suture Removal:} Non-dissolvable sutures may be removed at follow-up visits.
  \item \textbf{Imaging:} OCT may be performed to evaluate macular anatomy and detect complications.
  \item \textbf{Refractive Assessment:} A new glasses prescription may be needed after stabilization.
  \item \textbf{Secondary Procedures:} Cataract surgery or silicone oil removal may be necessary.
  \item \textbf{Warning Signs:} Patients are educated on warning signs that require immediate medical attention.
\end{itemize}

Adherence to the follow-up schedule and instructions is paramount for achieving the best possible long-term visual outcome.

\section*{Outcomes \& Prognosis}
Retinal detachment is a relatively uncommon but sight-threatening condition, with an estimated incidence of approximately 1 in 10,000 individuals per year for rhegmatogenous retinal detachment (RRD). It is more prevalent in older adults, particularly those in their 6th to 7th decades of life, although it can affect individuals of any age. High myopia is a significant risk factor, increasing the lifetime risk of detachment. Other risk factors include previous cataract surgery, ocular trauma, and a family history of retinal detachment. Males have a slightly higher incidence than females. Tractional retinal detachments are primarily associated with proliferative diabetic retinopathy, while exudative detachments are less common and typically secondary to underlying conditions.

The outcomes following retinal detachment surgery are generally favorable, with primary anatomical success rates ranging from 85\% to 95\%. However, anatomical success does not always guarantee full visual recovery. The most critical prognostic factor is the status of the macula at the time of surgery. If the macula is attached, the prognosis for good central vision is excellent. Conversely, if the macula is detached, the visual prognosis is guarded, with some degree of permanent vision loss often persisting.

Factors influencing prognosis include the duration of the detachment, the presence and severity of PVR, the size and location of retinal breaks, and the patient's overall ocular health. Recurrence of retinal detachment occurs in approximately 5-10\% of cases, often necessitating further surgical intervention. While the surgery aims to restore vision, patients should understand that final visual acuity may not return to pre-detachment levels, especially in complex cases.

\section*{References}
\begin{enumerate}
  \item MedlinePlus. Retinal detachment repair. U.S. National Library of Medicine; 2024. Available from: \url{https://medlineplus.gov/ency/article/002953.htm}
  
  \item American Academy of Ophthalmology. Preferred Practice Pattern: Retinal Detachment. San Francisco, CA: American Academy of Ophthalmology; 2023.
  
  \item Yanoff M, Duker JS. Ophthalmology. 6th ed. Elsevier; 2023.
  
  \item Ryan SJ, Sadda SVR, Hinton DR, et al. Ryan's Retina. 7th ed. Elsevier; 2022.
\end{enumerate}

\clearpage
